\emph{(i).} With $D = 1 + u + u^{2}$, $\sum_{i}p_{i} = (-u + 1 + u + u + u^{2})/D = 1$ and $\sum_{i}p_{i}^{2} = (u^{2} + (1+u)^{2} + u^{2}(1+u)^{2})/D^{2} = 1$, since $u^{2} + (1+u)^{2} + u^{2}(1+u)^{2} = (1 + u + u^{2})^{2}$; the script checks both as identities of rational functions. On $1 \le u \le 2$: $|p_{1}| = u/D$ has derivative $(1 - u^{2})/D^{2} \le 0$, so $|p_{1}| \ge 2/7$, with equality only at $u = 2$; $p_{2} = (1 + u)/D$ has derivative $-u(2 + u)/D^{2} < 0$, so $p_{2} \ge \tfrac37$; and $p_{3} = u(1 + u)/D$ has derivative $(1 + 2u)/D^{2} > 0$, so $p_{3} \ge \tfrac23$. That every era of the BKL map ends in an epoch with $1 \le u \le 2$ (the epoch map lowers $u$ by one until it enters $[1, 2)$, or reaches $1$ for rational data) is part of the bounce lemma filed alongside, synthesis.bkl-bounces-no-framewise-invariant; the classical BKL map is cited from memory, and the store's own Kasner data are the Schwarzschild exponents $p(1) = (-\tfrac13, \tfrac23, \tfrac23)$ of Theorem \claimref{interior.weighted-limit-theorem}.

\emph{(ii).} The proof of Theorem \claimref{interior.transport-level-convergence-along-tangential-curves} works at a single fixed time $\tau$ and uses three facts: the variation-of-constants identity, the bound $\|G - H\| \le \|E_{\tau}\|_{K}$, and the bound $\|H\| \le 1$ on unit directions, the last from Cauchy-Schwarz with $\sum_{i}p_{i}^{2} = 1$ and $\|RB_{i}R^{-1}\| = \|B_{i}\| = 1$ (conjugation by a rotation, an isometry for the norm used); on a unit direction $v$ the rotated model's generator is $-\sum_{i}(R\hat\theta)^{i}(v)\,p_{i}\,RB_{i}R^{-1}$, and $\bigl((R\hat\theta)^{i}(v)\bigr)_{i}$ is again a unit vector, so the same Cauchy-Schwarz bound applies. None of the three refers to a fixed Kasner point or frame, so the bound holds at time $\tau$ with any Kasner point and any rotation, in particular with the current epoch's model, and its constant $e^{\lambda(1 + \|E_{\tau}\|_{K})}$ is the same for all of them because $\sum p_{i}^{2} = 1$ on the whole circle, which every bounce maps to itself.

\emph{(iii).} The pairs $(u_{k}, R_{k})$ lie in the compact set $[1, 2] \times SO(3)$, so a subsequence converges to some $(u, R)$; along it the models converge, and since $\|E_{\tau_{k}}\|_{K} \to 0$ the values of $\widehat\Omega_{\tau_{k}}$ on the frame directions $R_{k}e_{i}$ converge to $-p_{i}(u)\,RB_{i}R^{-1}$ (the coframe $(R_{k}\hat\theta)^{j}$ takes the value $\delta^{j}_{i}$ on $R_{k}e_{i}$): these are accumulation values. By (i) all three $p_{i}(u)$ are nonzero, so each $RB_{i}R^{-1}$ lies in the generated algebra, and these three boosts generate $R\,\Lambda^{2}W\,R^{-1} = \Lambda^{2}W$ by Lemma \claimref{interior.boosts-generate-exactly-levi-factor} (conjugation by $R$ is a Lie algebra automorphism of $\Lambda^{2}W$). Conversely every accumulation value lies in $\Lambda^{2}W$: the Levi-Civita part of the Hubble-normalised connection takes values in $\Lambda^{2}W$ in a parallel orthonormal frame, and the tractor couplings on unit directions are $\tau$ times $\tau$-independent coefficients by the Einstein condition, the argument of clause (1) of Theorem \claimref{interior.weighted-limit-theorem}, so they tend to zero. Hence the generated algebra is exactly $\Lambda^{2}W$. The script computes the generated algebras in $\Cl(1,3)$: dimension $6$ from $-p_{i}B_{i}$ at $u = 1, \tfrac32, 2, \tfrac{11}{4}$, with two controls, dimension $3$ from two boosts and dimension $1$ from the single boost of the Taub point.

\emph{Scope.} The hypothesis of (iii) is that the approach returns arbitrarily close to non-Taub Kasner models; for a BKL approach with irrational parameter the last epoch of each era provides such returns in the BKL picture, and for Bianchi IX vacuum Ringstr\"om's attractor theorem gives convergence of the Hubble-normalised state to the attractor containing the Kasner circle (both cited from memory, not proved here and not store entities). Rational parameters end at $u = 1$ and then reach a Taub point (the flat Kasner point), where (iii) does not apply.
